\documentclass[../../../include/open-logic-section]{subfiles}

\begin{document}

\olfileid{his}{set}{cantorplane}
\olsection{રેખા અને સમતલ વિશે કૅન્ટૉર}

શ્રેડર--બર્નસ્ટાઇન પ્રમેયની સાબિતીની
આસપાસની કેટલીક ઘટનાઓ
\olref[his][set][pathology]{sec}માં
ચર્ચેલા ઇતિહાસ સાથે જોડાય છે.
યાદ કરો કે 1877માં કૅન્ટૉરે
ચોરસમાં અને તેની એક બાજુમાં
બિંદુઓની સંખ્યા સરખી છે
તે સાબિત કર્યું હતું. અહીં
તેમની પહેલી પ્રયત્નરૂપ સાબિતી રજૂ કરીએ.

$\unitline$ એકમ રેખા, એટલે
બિંદુઓનો ગણ $[0,1]$,
અને $\unitsquare$ એકમ ચોરસ,
એટલે બિંદુઓનો ગણ $\unitline
\times \unitline$, માનો.
આ સંકેતમાં કૅન્ટૉરે
$\cardeq{\unitline}{\unitsquare}$
સાબિત કર્યું. ડેડેકિન્ડને
લખેલી નોંધમાં તેમણે
મોટા ભાગે નીચેની દલીલ આપી.

\begin{thm}\ollabel{thm:cantorplane}
$\cardeq{\unitline}{\unitsquare}$
\end{thm}

\begin{proof}[સાબિતી: પ્રથમ ભાગ.]
$a, b \in \unitline$ નિશ્ચિત કરો.
તેમને દ્વિઆધારી સંકેતમાં લખો,
જેથી $0$ અને $1$ના અનંત
અનુક્રમો $a_1$, $a_2$, \dots,
અને $b_1$, $b_2$, \dots,
આ રીતે મળે:
\begin{align*}
a &= 0.a_1a_2a_3a_4\dots\\
b &= 0.b_1b_2b_3b_4\dots
\intertext{હવે આ રીતે આપેલું વિધેય $f \colon \unitsquare \to \unitline$ લો:}
f(a, b) & = 0.a_1b_1a_2b_2a_3b_3a_4b_4\dots
\end{align*}
હવે $f$ એ !!a{injection} છે,
કારણ કે જો $f(a, b) = f(c,d)$,
તો દરેક $n \in \Nat$ માટે
$a_n = c_n$ અને $b_n = d_n$;
તેથી $a = c$ અને $b = d$.
\end{proof}

દુર્ભાગ્યે, ડેડેકિન્ડે કૅન્ટૉરને
દર્શાવ્યું તેમ, આ મૂળ પ્રશ્નનો
ઉત્તર આપતું નથી. $0.\dot{1}\dot{0} =
0.1010101010\ldots$ લો.
આને $f(a,b) = 0.\dot{1}\dot{0}$
થવા માટે નીચેના અંકો જોઈએ:
\begin{align*}
a&= 0.\dot{1}\dot{1} = 0.111111\ldots\\
b&= 0
\end{align*}
પરંતુ $a = 0.\dot{1}\dot{1} = 1$ છે.
એટલે “$a$ અને $b$ને દ્વિઆધારી
સંકેતમાં લખો” કહીએ ત્યારે
\emph{કયો} સંકેત લેવો તે પસંદ કરવું
પડે; અને $f$ \emph{વિધેય} હોવું
જોઈએ એટલે બે સંભવિત
સંકેતોમાંથી ફક્ત \emph{એક}
જ લઈ શકીએ. દાખલા તરીકે,
સરળ સંકેત લઈ $a$ને
“$1.000\ldots$” લખીએ,
તો $f(a, b) = 0.\dot{1}\dot{0}$
થાય એવી કોઈ જોડી
$\tuple{a, b}$ મળતી નથી.
\sourcecorrection{OLMD-169}{મૂળનું આ
ઉદાહરણ પસંદ કરેલા સંકેત પર
આધારિત છે: $1$ને
$0.111\ldots$ લખીએ તો
$0.101010\ldots$ ખરેખર
$f(1,0)$ બને છે; અને
$1.000\ldots$ એ ઉપરના
$0.a_1a_2\ldots$ સ્વરૂપમાં નથી.
અવ્યાપ્તતાનું માન્ય ઉદાહરણ
માટે $1$ સિવાયની સંખ્યાઓનું
અંતે માત્ર $1$ આવતું ન હોય
એવું દ્વિઆધારી સ્વરૂપ,
અને $1=0.111\ldots$,
એકસરખી રીતે પસંદ કરો.
પછી $0.00101010\ldots$
સહપ્રદેશમાં છે, પણ તેના
વિષમ સ્થાનોના અંકો
$0.01111\ldots=\frac12$
બને છે; $\frac12$ માટે
પસંદ થયેલું સ્વરૂપ
$0.1000\ldots$ હોવાથી
આ સંખ્યા $f$ની કિંમત
નથી.}

સારાંશમાં, ડેડેકિન્ડે દર્શાવ્યું
કે કેટલાક પુનરાવર્તિત
દ્વિઆધારી વિસ્તરણો શક્ય હોવાથી
કૅન્ટૉરનું વિધેય $f$
!!a{injection} છે, પણ
!!a{surjection} \emph{નથી}.
તેથી કૅન્ટૉરે ફક્ત
$\cardle{\unitsquare}{\unitline}$
બતાવ્યું છે,
$\cardeq{\unitsquare}{\unitline}$
\emph{નહીં}.

કૅન્ટૉરે લગભગ તરત જ
ડેડેકિન્ડને જવાબ લખીને
સાબિતી નીચે મુજબ
પૂરી કરી શકાય એવું સૂચવ્યું:

\begin{proof}[સાબિતી: પૂર્ણ ભાગ.]
આપણે $\cardle{\unitsquare}{\unitline}$
બતાવ્યું છે. પણ $\unitline$માંથી
$\unitsquare$માં !!a{injection}
સ્પષ્ટ છે: ચોરસની એક બાજુ
પર રેખાને મૂકી દો. તેથી
$\cardle{\unitline}{\unitsquare}$ અને
$\cardle{\unitsquare}{\unitline}$. શ્રેડર--બર્નસ્ટાઇન પ્રમેયથી
(\olref[sfr][siz][sb]{thm:schroder-bernstein}),
$\cardeq{\unitline}{\unitsquare}$.
\end{proof}

અલબત્ત, કૅન્ટૉર આ શબ્દોમાં
છેલ્લી પંક્તિ પૂરી કરી શકતા
નહોતા, કારણ કે શ્રેડર--બર્નસ્ટાઇન
પ્રમેય ત્યારે હજી સાબિત થયું
નહોતું. કૅન્ટૉરે પછી તેને
સામાન્ય અનુમાન તરીકે રજૂ કર્યું,
પરંતુ 1897 સુધી તેની
સંતોષકારક સાબિતી મળી નહોતી.
(એટલે 1877માં જ પછીથી
કૅન્ટૉરે \olref{thm:cantorplane}
માટે શ્રેડર--બર્નસ્ટાઇનનો
આધાર ન લેતી બીજી
સાબિતી આપી.)

\end{document}
